\uuid{C38u}
\titre{Base et dimension}
\auteur{hanani ; mzari}
\organisation{exo7}
\licence{CC BY-NC-SA 4.0}
\datecreate{2020-01-20}
\niveau{L1}
\module{Algèbre}
\chapitre{Espace vectoriel}
\sousChapitre{Base}
\difficulte{3}
\theme{Base et dimension}
\exo7qcmid{lille-2.espaces-vectoriels.047}

\contenu{
\begin{question}
\qclassification{106.04}
\qtags{motcle=Bases de $\Cc^2$}
Quelles sont les assertions vraies ?
\begin{answers}
\good{$\{(1,0),(1,1)\}$ est une base de $\Cc^2$ comme $\Cc$-espace vectoriel}.
\good{La dimension de $\Cc^2$ comme $\Rr$-espace vectoriel est $4$}.
\good{$\{(1,0),(0,i),(i,0),(0,1)\}$ est une base de $\Cc^2$ comme $\Rr$-espace vectoriel}.
\bad{La dimension de $\Cc^2$ comme $\Rr$-espace vectoriel est $2$}.
\end{answers}
\begin{explanations} L'espace $\Cc^2$ est un $\Cc$-espace vectoriel de dimension $2$. Par conséquent, pour tous $(a,b), (c,d) \in \Cc^2$, non colinéaires sur $\Cc$, $\{(a,b), (c,d)\}$ est une $\Cc$-base de $\Cc^2$.
\vskip2mm
D'autre part $\Cc^2$ est un  $\Rr$-espace vectoriel de dimension $4$. Par conséquent, toute famille $\{(a,b), (a',b'), (c,d), (c',d')\}$, de vecteurs de $\Cc^2$ linéairement indépendants sur $\Rr$, est une $\Rr$-base de $\Cc^2$.
\end{explanations}
\end{question}
}
